Every Lightning implementation disclosed security vulnerabilities in 2026, and every advisory told operators to
upgrade. Whether they do has never been measured, because Lightning gossip carries no software version. We show that
the release line of a public node can be read from the feature bits it announces: matched against the source of every
release of LND, Core Lightning, Eclair and LDK, the inferred range contains the true line on \PctVerStrict\% of
\NVerNodeDays{} node-days of nodes that state their version, and the most likely line is exact on \PctVerLikely\%.
Applying it to \DataUsableDays{} daily snapshots of the public graph from 2023 to 2026, we find that half of the
LND nodes on an older line upgrade within \MedianLNDAll{} days of a release, in waves consistent with node platforms
pushing updates, while small nodes and nodes that skipped earlier releases lag far behind. One operator ran up to
\FleetPeakNodes{} public LND nodes, \FleetPeakShare\% of all live LND nodes, with about one bitcoin between them.
For \NAdvTotal{} LND and CLN advisories, fixes shipped a median \FixToDiscMedian{} days before disclosure. LND's
disclosures did not change how fast operators left vulnerable releases: two months after Lightning Labs' August~2026
batch, \PctCohortStillLLBatch\% of the nodes exposed at disclosure still ran an affected release. CLN's disclosures
of July~2026 raised the monthly rate at which nodes left vulnerable releases two- to fourfold. Exposure measured by
capacity is small; by node count, it depends on how one counts a single operator's fleet.
